\documentclass[a4paper,11pt]{article}

% ============================================================
%  Le Collé — Séance 1 : Déterminants, courbes paramétrées
%  © 2026 Ahmed Amine Ouqlani
% ============================================================

% ---------- Polices ----------
\usepackage[utf8]{inputenc}
\usepackage[T1]{fontenc}
\usepackage[sc]{mathpazo}          % texte et mathématiques : Palatino
\usepackage[scaled=0.88]{helvet}   % titres et repères : linéale
\linespread{1.07}
\usepackage{amsmath,amssymb,amsthm,mathrsfs,stmaryrd}
\usepackage{microtype}

% ---------- Mise en page ----------
\usepackage[a4paper,top=2.3cm,bottom=2.3cm,left=2.4cm,right=2.4cm,
            headheight=14pt,headsep=0.8cm,footskip=1.1cm]{geometry}
\setlength{\parindent}{0pt}
\setlength{\parskip}{0.45em plus 0.15em minus 0.1em}
\raggedbottom
\emergencystretch=1.5em
\clubpenalty=10000 \widowpenalty=10000 \displaywidowpenalty=10000 \brokenpenalty=10000

% ---------- Couleurs : un bleu de nuit, un ocre, des gris ----------
\usepackage{xcolor}
\definecolor{Ink}{HTML}{1C2533}
\definecolor{Bleu}{HTML}{1F3A5F}
\definecolor{BleuMoyen}{HTML}{4A6A8F}
\definecolor{BleuClair}{HTML}{9DB2C9}
\definecolor{Fond}{HTML}{F2F4F7}
\definecolor{Ocre}{HTML}{B07A26}
\definecolor{Brique}{HTML}{A23B2A}
\definecolor{Gris}{HTML}{5E6672}
\definecolor{Filet}{HTML}{C5CCD6}

% ---------- Packages ----------
\usepackage{enumitem}
\setlist{itemsep=0.25em,topsep=0.35em,parsep=0pt}
\setlist[itemize]{label=\textcolor{BleuMoyen}{\small\textbullet},leftmargin=1.4em}
\setlist[enumerate]{leftmargin=1.8em,label=\textsf{\small\color{Bleu}\arabic*.}}
\usepackage{booktabs,array,multicol}
\usepackage{needspace,float}
\usepackage[font=small,labelfont={sf,color=Bleu},labelsep=period,skip=6pt]{caption}
\usepackage{tikz}
\usetikzlibrary{arrows.meta,calc,decorations.markings}
\usepackage{pgfplots}
\pgfplotsset{compat=1.17}
\usepackage{tkz-tab}
\usepackage[most]{tcolorbox}
\usepackage{titlesec,titletoc}
\usepackage{fancyhdr,lastpage,etoolbox}
\usepackage[colorlinks=true,linkcolor=Bleu,urlcolor=Bleu,citecolor=Bleu,
            pdftitle={Le Collé — TD Réduction (énoncés)},pdfauthor={Ahmed Amine Ouqlani}]{hyperref}

\numberwithin{equation}{section}
\numberwithin{figure}{section}
\numberwithin{table}{section}

% ---------- Titres ----------
% Chapitres (sections) : grand numéro, titre en linéale, filet ocre
\titleformat{\section}[display]
  {\normalfont\sffamily\color{Ink}}
  {{\small\color{Gris}\textls[200]{\MakeUppercase{Chapitre}}}\\[-0.2em]%
   {\fontsize{64}{64}\selectfont\color{BleuClair}\thesection}}
  {0.4em}
  {\fontsize{25}{30}\selectfont\bfseries}
  [\vspace{0.5em}{\color{Ocre}\rule{3.2cm}{1.4pt}}]
\titlespacing*{\section}{0pt}{0pt}{1.6em}
\titleformat{name=\section,numberless}[display]
  {\normalfont\sffamily\color{Ink}}{}{0pt}{\fontsize{25}{30}\selectfont\bfseries}
  [\vspace{0.5em}{\color{Ocre}\rule{3.2cm}{1.4pt}}]
\titlespacing*{name=\section,numberless}{0pt}{0pt}{1.6em}

\titleformat{\subsection}
  {\needspace{5\baselineskip}\normalfont\sffamily\Large\bfseries\color{Ink}}
  {\color{Bleu}\thesubsection}{0.8em}{}
\titlespacing*{\subsection}{0pt}{1.8em}{0.5em}
\titleformat{\subsubsection}
  {\needspace{5\baselineskip}\normalfont\sffamily\large\bfseries\color{Ink}}
  {\color{Bleu}\thesubsubsection}{0.7em}{}
\titlespacing*{\subsubsection}{0pt}{1.6em}{0.4em}
\titleformat{\paragraph}[runin]{\normalfont\sffamily\bfseries\color{Ink}}{}{0pt}{}[.]
\titlespacing*{\paragraph}{0pt}{0.9em}{0.6em}
\setcounter{secnumdepth}{3}
\setcounter{tocdepth}{2}

% ---------- Table des matières ----------
\titlecontents{section}[0pt]{\addvspace{0.7em}\sffamily\bfseries\color{Ink}}
  {\makebox[2.2em][l]{\color{Bleu}\thecontentslabel}}{}{\hfill\contentspage}[\addvspace{0.2em}]
\titlecontents{subsection}[2.2em]{\small\linespread{1}\selectfont}
  {\makebox[2.6em][l]{\color{Gris}\thecontentslabel}}{}{\titlerule*[0.7pc]{.}\contentspage}

% ---------- En-têtes et pieds de page ----------
\pagestyle{fancy}
\fancyhf{}
\renewcommand{\sectionmark}[1]{\markboth{\thesection\quad#1}{}}
\fancyhead[L]{\footnotesize\sffamily\color{Gris}\leftmark}
\fancyhead[R]{\footnotesize\sffamily\scshape\color{Gris}Le Collé}
\fancyfoot[C]{\small\sffamily\color{Ink}\thepage}
\fancyfoot[L]{\scriptsize\sffamily\hyperlink{toc}{\color{Filet}sommaire}}
\renewcommand{\headrulewidth}{0.4pt}
\renewcommand{\headrule}{{\color{Filet}\hrule height\headrulewidth width\headwidth}}
\renewcommand{\footrulewidth}{0pt}
\fancypagestyle{plain}{\fancyhf{}\fancyfoot[C]{\small\sffamily\color{Ink}\thepage}\renewcommand{\headrulewidth}{0pt}}

% ---------- Énoncés ----------
\newtheoremstyle{enonce}{0pt}{0pt}{}{}{\sffamily\bfseries\color{Bleu}}{.}{0.5em}
  {\thmname{#1}\thmnumber{~#2}\thmnote{\ \normalfont\sffamily\color{Ink}(#3)}}
\newtheoremstyle{exemple}{0.9em}{0.9em}{}{}{\sffamily\bfseries\color{BleuMoyen}}{.}{0.5em}
  {\thmname{#1}\thmnumber{~#2}\thmnote{\ \normalfont\sffamily\color{Ink}(#3)}}
\newtheoremstyle{remarque}{0.7em}{0.7em}{}{}{\sffamily\itshape}{.}{0.5em}{\thmname{#1}\thmnote{\ (#3)}}
\theoremstyle{enonce}
\newtheorem{definition}{Définition}[section]
\newtheorem{theoreme}[definition]{Théorème}
\newtheorem{proposition}[definition]{Proposition}
\newtheorem{corollaire}[definition]{Corollaire}
\theoremstyle{exemple}
\newtheorem{exemple}[definition]{Exemple}
\theoremstyle{remarque}
\newtheorem*{remarque}{Remarque}

\tcbset{
  defi/.style={enhanced,breakable,colback=white,colframe=white,boxrule=0pt,arc=0pt,
     borderline west={1.2pt}{0pt}{Bleu},left=11pt,right=4pt,top=3pt,bottom=3pt,
     before skip=1.1em,after skip=1.1em},
  thm/.style={enhanced,breakable,colback=Fond,colframe=Fond,boxrule=0pt,arc=0pt,
     borderline west={2.4pt}{0pt}{Bleu},left=11pt,right=8pt,top=6pt,bottom=6pt,
     before skip=1.2em,after skip=1.2em}}
\tcolorboxenvironment{definition}{defi}
\tcolorboxenvironment{theoreme}{thm}
\tcolorboxenvironment{proposition}{thm}
\tcolorboxenvironment{corollaire}{thm}

\renewcommand{\proofname}{Démonstration}
\renewcommand{\qedsymbol}{\textcolor{BleuClair}{$\blacksquare$}}
\makeatletter
\renewenvironment{proof}[1][\proofname]{\par\pushQED{\qed}\normalfont
  \topsep0.5em\relax\trivlist\item[\hskip\labelsep\sffamily\itshape\color{Gris}#1.]\ignorespaces}
  {\popQED\endtrivlist\@endpefalse}
\makeatother

% ---------- Méthode : filets haut et bas, titre en bandeau ----------
\makeatletter\newcommand{\methlabelfix}{\protected@edef\@currentlabel{\thetcbcounter}}\makeatother
\newtcolorbox[auto counter,number within=section]{methode}[2][]{%
  enhanced,unbreakable,colback=white,colframe=white,boxrule=0pt,arc=0pt,
  borderline north={0.9pt}{0pt}{Bleu},borderline south={0.9pt}{0pt}{Bleu},
  left=4pt,right=4pt,top=8pt,bottom=8pt,before skip=1.4em,after skip=1.4em,
  before upper={\methlabelfix\phantomsection\addcontentsline{met}{metentry}{\protect\makebox[2.6em][l]{\thetcbcounter}#2}%
     {\sffamily\bfseries\small\color{white}\colorbox{Bleu}{\strut\ MÉTHODE~\thetcbcounter\ }}%
     \ \ {\sffamily\bfseries\color{Bleu}#2}\par\smallskip},#1}

% ---------- Repères pédagogiques : un filet de couleur, un mot-clé ----------
\newcommand{\definirbloc}[3]{% nom, couleur, libellé
  \newtcolorbox{#1}[1][]{enhanced,breakable,blanker,
    borderline west={1.4pt}{0pt}{#2},left=10pt,right=2pt,top=1pt,bottom=1pt,
    before skip=1em,after skip=1em,
    before upper={{\sffamily\bfseries\footnotesize\color{#2}\textls[120]{\MakeUppercase{#3}}}\quad},##1}}
\definirbloc{reflexe}{Bleu}{Réflexe}
\definirbloc{astuce}{Ocre}{Astuce}
\definirbloc{attention}{Brique}{Attention}
\definirbloc{enpratique}{BleuMoyen}{En pratique}
\definirbloc{plusloin}{Gris}{Pour aller plus loin}
\definirbloc{strategie}{Bleu}{Stratégie}
\definirbloc{lecture}{BleuMoyen}{Lire le schéma}
% compatibilité avec les parties reprises
\newenvironment{garde}{\begin{attention}}{\end{attention}}
\newenvironment{complementHP}{\begin{plusloin}}{\end{plusloin}}
\newenvironment{interpretation}{\begin{enpratique}}{\end{enpratique}}

% ---------- Exercices, corrigés ----------
\newcounter{exo}[section]
\renewcommand{\theexo}{\thesection.\arabic{exo}}
\newcommand{\niveau}[1]{{\color{Ocre}\footnotesize
  \ifcase#1\or$\bullet\circ\circ\circ\circ$\or$\bullet\bullet\circ\circ\circ$\or$\bullet\bullet\bullet\circ\circ$%
  \or$\bullet\bullet\bullet\bullet\circ$\or$\bullet\bullet\bullet\bullet\bullet$\fi}}
\newcommand{\serie}[1]{\par\needspace{8\baselineskip}\vspace{1.5em}%
  \noindent{\sffamily\bfseries\color{Bleu}#1}\par\nobreak\vspace{-0.4em}%
  \noindent{\color{Filet}\rule{\linewidth}{0.4pt}}\par\nobreak\vspace{0.1em}}
\newenvironment{exo}[3]{%
  \par\needspace{5\baselineskip}\vspace{0.75em}\refstepcounter{exo}\label{e:#1}%
  \def\exolab{#1}%
  \noindent{\sffamily\bfseries\color{Bleu}Exercice~\theexo}\quad\niveau{#2}%
  \hfill{\footnotesize\sffamily\color{Gris}#3}\par\nobreak\vspace{0.15em}}%
  {\par\nopagebreak\noindent\hfill{\scriptsize\sffamily\hyperref[c:\exolab]{\color{Gris}corrigé p.~\pageref{c:\exolab}}}\par}
\newenvironment{corrige}[1]{%
  \par\needspace{6\baselineskip}\vspace{1em}\phantomsection\label{c:#1}%
  \noindent{\sffamily\bfseries\color{Bleu}Exercice~\ref{e:#1}}\quad
  {\scriptsize\sffamily\hyperref[e:#1]{\color{Gris}énoncé p.~\pageref{e:#1}}}\par\nobreak\vspace{0.2em}}%
  {\hfill\textcolor{BleuClair}{$\square$}\par}
\newcommand{\etape}[1]{\par\smallskip\noindent{\sffamily\bfseries\small\color{BleuMoyen}#1.}\ }
\newenvironment{solution}{\par\noindent{\sffamily\bfseries\small\color{BleuMoyen}Solution.}\ }{\hfill\textcolor{BleuClair}{$\square$}\par}
\newcommand{\chercher}{\par\smallskip\noindent{\color{Filet}\rule[0.5ex]{1.2cm}{0.4pt}}\
  {\footnotesize\sffamily\itshape\color{Gris}laisser chercher l'étudiant avant de poursuivre}\par\smallskip}

% ---------- Guide de séance ----------

\newcommand{\marge}[1]{\leavevmode\marginpar{\raggedleft\scriptsize\sffamily\color{Gris}\textls[100]{\MakeUppercase{#1}}}\ignorespaces}
\newcommand{\Etape}[4]{%
  \subsection[Étape #1 --- #3 {\color{Gris}(#2)}]{#3}\label{etape:#1}%
  \noindent{\small\sffamily\color{Gris}{\color{Ocre}\bfseries Étape #1}\quad·\quad #2\quad·\quad #4}\par\medskip}
\newcommand{\question}[1]{\par\smallskip\noindent{\sffamily\bfseries\small\color{BleuMoyen}Question.}\ \textit{#1}\par\smallskip}
\newcommand{\rappel}[1]{\par\smallskip\noindent\phantomsection\addcontentsline{rap}{rapentry}{#1}%
  {\sffamily\bfseries\scriptsize\colorbox{Fond}{\color{Gris}\strut RAPPEL}}\ {\sffamily\small\color{Gris}#1.}\ }
\newcommand{\pourquoi}{\par\noindent{\sffamily\bfseries\small\color{Ocre}Pourquoi ?}\ }

% Ouverture de chapitre : \ouverture{objectifs}{prérequis}{durée}
\newcommand{\ouverture}[3]{%
  \begin{tcolorbox}[enhanced,colback=Fond,colframe=Fond,arc=0pt,boxrule=0pt,left=12pt,right=12pt,top=10pt,bottom=10pt,
     before skip=1.2em,after skip=1.6em]
  \begin{minipage}[t]{0.56\linewidth}
  {\sffamily\bfseries\small\color{Bleu}OBJECTIFS}\par\vspace{0.3em}\small #1
  \end{minipage}\hfill
  \begin{minipage}[t]{0.38\linewidth}
  {\sffamily\bfseries\small\color{Bleu}PRÉREQUIS}\par\vspace{0.3em}\small #2
  \ifx\relax#3\relax\else\par\vspace{0.6em}{\sffamily\bfseries\small\color{Bleu}DURÉE}\par\vspace{0.2em}\small #3\fi
  \end{minipage}
  \end{tcolorbox}}
% Fin de chapitre
\newenvironment{essentiel}{\par\needspace{12\baselineskip}\vspace{1.5em}%
  \begin{tcolorbox}[enhanced,unbreakable,colback=white,colframe=white,boxrule=0pt,arc=0pt,
   borderline north={1.4pt}{0pt}{Ocre},left=2pt,right=2pt,top=8pt,bottom=4pt]
  {\sffamily\bfseries\color{Ocre}L'ESSENTIEL DU CHAPITRE}\par\medskip}{\end{tcolorbox}}

% ---------- Notations ----------
\newcommand{\K}{\mathbb{K}}
\newcommand{\R}{\mathbb{R}}
\newcommand{\C}{\mathbb{C}}
\newcommand{\N}{\mathbb{N}}
\newcommand{\Z}{\mathbb{Z}}
\newcommand{\Mn}[1]{\mathcal{M}_{n}(#1)}
\newcommand{\GLn}[1]{\mathrm{GL}_n(#1)}
\newcommand{\LE}{\mathcal{L}(E)}
\newcommand{\rg}{\operatorname{rg}}
\newcommand{\vect}{\operatorname{Vect}}
\newcommand{\Mat}{\operatorname{Mat}}
\newcommand{\Id}{\operatorname{Id}}
\newcommand{\sh}{\operatorname{sh}}
\newcommand{\ch}{\operatorname{ch}}
\renewcommand{\th}{\operatorname{th}}
\DeclareMathOperator{\argsh}{argsh}
\newcommand{\vv}[1]{\overrightarrow{#1}}
\newcommand{\ve}[1]{\vec{#1}}
\newcommand{\ient}[2]{\llbracket #1,#2\rrbracket}
\newcommand{\dt}[1]{\left|\begin{matrix}#1\end{matrix}\right|}
\newcommand{\dis}{\displaystyle}

% ---------- Figures ----------
\pgfplotsset{trig format plots=rad}
\pgfplotsset{courbe/.style={axis lines=middle,axis equal image,
  axis line style={Gris,-{Stealth[length=5pt]}},
  xlabel={$x$},ylabel={$y$},ticklabel style={font=\scriptsize\sffamily,color=Gris},
  every axis x label/.style={at={(ticklabel* cs:1)},anchor=west,font=\small},
  every axis y label/.style={at={(ticklabel* cs:1)},anchor=south,font=\small},
  samples=300,smooth,width=0.6\linewidth}}
\tikzset{arcetude/.style={Bleu,line width=1.25pt},
         reste/.style={BleuClair,line width=1pt},
         tgte/.style={Ocre,line width=0.9pt},
         asy/.style={Gris,dashed,line width=0.6pt},
         pt/.style={circle,fill=Ink,inner sep=1.2pt}}

% ---------- Navigation ----------
\makeatletter
\newcommand*\l@rapentry{\@dottedtocline{1}{0em}{0em}}
\newcommand*\l@metentry{\@dottedtocline{1}{0em}{0em}}
\newcommand{\listemethodes}{\@starttoc{met}}
\newcommand{\listerappels}{\@starttoc{rap}}
\makeatother
% onglets de chapitre sur le bord droit
\usepackage{eso-pic}
\newif\ifonglets \ongletsfalse
\newcommand{\dessineronglet}{%
  \edef\ongcol{\ifnum\value{section}=3 Ocre\else Bleu\fi}%
  \edef\ongnum{\the\value{section}}%
  \pgfmathsetmacro{\ongy}{3.2+2.3*(\ongnum-1)}%
  \begin{tikzpicture}[remember picture,overlay]
  \fill[\ongcol] ([xshift=-0.75cm,yshift=-\ongy cm]current page.north east)
     rectangle ([yshift=-\ongy cm-1.9cm]current page.north east);
  \node[text=white,font=\sffamily\bfseries\large] at ([xshift=-0.37cm,yshift=-\ongy cm-0.95cm]current page.north east)
     {\ongnum};
  \end{tikzpicture}}
\AddToShipoutPictureBG{\ifonglets\ifnum\value{section}>0 \dessineronglet\fi\fi}
% blocs de navigation et d'enseignement
\definirbloc{commencer}{Gris}{Si vous commencez ici}
\definirbloc{horsprogramme}{Brique}{Hors programme}
\newcommand{\HP}{\,{\sffamily\bfseries\scriptsize\textcolor{Brique}{[HP]}}}
\newcommand{\source}[1]{{\footnotesize\sffamily\color{Gris}\,(#1)}}
\newcommand{\coupdepouce}[1]{\par\vspace{-0.2em}\noindent{\footnotesize\sffamily\color{Gris}%
  \textbf{Coup de pouce} (à donner si l'étudiant bloque) : #1}\par}
\newcommand{\pointarret}[1]{\par\vspace{1em}\noindent{\color{Filet}\rule[0.5ex]{\linewidth}{0.4pt}}\par\nobreak
  \noindent{\footnotesize\sffamily\color{Ocre}\bfseries POINT D'ARRÊT POSSIBLE}\quad{\footnotesize\sffamily\color{Gris}#1}\par\vspace{0.6em}}
\newcommand{\demarche}[1]{\noindent{\sffamily\bfseries\small\color{BleuMoyen}Démarche.}\ \textit{#1}\par\smallskip}

% ---------- Corrigés : idées, astuces, pièges, à retenir ----------
\newcommand{\choix}[1]{\par\noindent{\sffamily\bfseries\small\color{BleuMoyen}Choix de la méthode.}\ #1\par\smallskip}
\newcommand{\gain}[1]{\par\noindent{\sffamily\bfseries\small\color{Ocre}Astuce.}\ #1\par}
\newcommand{\piege}[1]{\par\noindent{\sffamily\bfseries\small\color{Brique}Piège.}\ #1\par}
\newtcolorbox{retenirbox}{enhanced,breakable,colback=Fond,colframe=Fond,arc=0pt,boxrule=0pt,
  borderline west={2pt}{0pt}{Ocre},left=8pt,right=6pt,top=4pt,bottom=4pt,before skip=0.7em,after skip=0.5em}
\newcommand{\retenir}[1]{\begin{retenirbox}{\sffamily\bfseries\footnotesize\color{Ocre}À RETENIR}\quad #1\end{retenirbox}}
\newcommand{\TD}[1]{{\color{Ocre}\bfseries TD\,#1}\ ·\ }
\newcommand{\q}[1]{\par\smallskip\noindent{\sffamily\bfseries\color{Bleu}#1}\ }
\providecommand{\iddots}{\mathinner{\mkern1mu\raise1pt\hbox{.}\mkern2mu\raise4pt\hbox{.}\mkern2mu\raise7pt\hbox{.}\mkern1mu}}

\providecommand{\Ker}{\operatorname{Ker}}\renewcommand{\Im}{\operatorname{Im}}\providecommand{\tr}{\operatorname{tr}}\providecommand{\Sp}{\operatorname{Sp}}

% ================================================================ TD : mise en page propre au livret d'exercices
\fancyhf{}
\fancyhead[L]{\footnotesize\sffamily\color{Gris}TD · Réduction · énoncés}
\fancyhead[R]{\footnotesize\sffamily\scshape\color{Gris}Le Collé}
\fancyfoot[C]{\small\sffamily\color{Ink}\thepage\,/\,\pageref{LastPage}}
\fancypagestyle{plain}{\fancyhf{}\fancyfoot[C]{\small\sffamily\color{Ink}\thepage\,/\,\pageref{LastPage}}\renewcommand{\headrulewidth}{0pt}}
\makeatletter
\newcommand{\partie}[3]{% numéro, titre, description
  \par\needspace{10\baselineskip}\vspace{1.4em}\phantomsection\label{partie:#1}\pdfbookmark[0]{Partie #1}{partie.#1}%
  {\sffamily\small\color{Gris}\textls[200]{PARTIE #1}}\par\nobreak\vspace{0.1em}
  {\sffamily\bfseries\LARGE\color{Ink}#2}\par\nobreak\vspace{0.3em}
  {\color{Ocre}\rule{2.6cm}{1.2pt}}\par\nobreak\vspace{0.2em}
  {\small\color{Gris}#3}\par\nobreak\vspace{0.6em}}
\newcommand{\exercice}[7]{% numéro, difficulté, titre, origine, notions, identifiant, page du corrigé
  \par\needspace{8\baselineskip}\vspace{1.1em}\phantomsection
  \def\@currentlabel{#1}\label{e:#6}\pdfbookmark[1]{Exercice #1}{e.#6}%
  \noindent{\sffamily\bfseries\large\color{Bleu}Exercice~#1}\quad\niveau{#2}\hfill{\footnotesize\sffamily\color{Gris}#4}\par\nobreak
  \vspace{0.1em}\noindent{\sffamily\bfseries\color{Ink}#3}\par\nobreak
  \vspace{0.1em}\noindent{\footnotesize\sffamily\color{Gris}Notions : #5.\hfill\mbox{corrigé p.~#7}\par}\nobreak
  \vspace{0.35em}}
\newcommand{\renvoi}[1]{\par\nopagebreak\vspace{0.2em}\noindent\hfill{\scriptsize\sffamily\color{Gris}#1}\par}
\newcommand{\place}{\par\vspace{1.6cm plus 0.6cm minus 0.6cm}}
\newcommand{\correction}[4]{% numéro, titre, identifiant, renvoi
  \par\needspace{7\baselineskip}\vspace{1.3em}\phantomsection
  \def\@currentlabel{#1}\label{e:#3}\pdfbookmark[1]{Exercice #1}{c.#3}%
  {\color{Filet}\hrule height 0.4pt}\vspace{0.5em}
  \noindent{\sffamily\bfseries\large\color{Bleu}Exercice~#1}\quad{\sffamily\bfseries\color{Ink}#2}\hfill{\scriptsize\sffamily\color{Gris}#4}\par\nobreak
  \vspace{0.4em}}
\newcommand{\sommaireligne}[3]{\noindent\hbox to 2.2em{\hfill\sffamily\color{Bleu}#1.}\enspace\begin{minipage}[t]{0.8\linewidth}\raggedright#2\end{minipage}\hfill\pageref{e:#3}\par}
\makeatother
\setlength{\parskip}{0.5em plus 0.15em minus 0.1em}

\begin{document}
\thispagestyle{plain}
{\sffamily\small\color{Gris}\textls[150]{LE COLLÉ · MATHÉMATIQUES PT · CHAPITRE 3}}\par\vspace{0.2em}
{\sffamily\bfseries\fontsize{26}{30}\selectfont\color{Ink}TD --- Réduction}\par\vspace{0.3em}
{\sffamily\Large\color{BleuMoyen}Énoncés}\par\vspace{0.4em}
{\color{Ocre}\rule{3.2cm}{1.4pt}}\par\vspace{0.8em}
{\small Ce livret rassemble 24 exercices, classés en quatre parties, du plus direct au plus exigeant.
Il reprend \textbf{tous les exercices du livret} (proposés à la fin de chaque partie du cours), complétés par des
\textbf{exercices complémentaires}, absents du livret. L'origine de chaque exercice est indiquée en haut à droite :
\emph{« Livret, exercice~1.5 »} ou \emph{« Exercice complémentaire »}. La difficulté va de \niveau{1} à \niveau{5}.
Les corrections sont dans un livret séparé, \texttt{TD\_red\_corrections.pdf} ; la page de chaque correction est indiquée en tête de l'énoncé.}

\vspace{0.6em}{\sffamily\bfseries\color{Bleu}Sommaire}\par\vspace{0.2em}{\small\setlength{\parskip}{1.5pt}\setlength{\columnsep}{1.2cm}\begin{multicols}{2}\raggedcolumns
\vspace{0.5em}\noindent\begin{minipage}[t]{0.86\linewidth}\raggedright\leavevmode\sffamily\bfseries\color{Ink}I. Applications directes du cours\end{minipage}\hfill\pageref{partie:I}\par\vspace{0.15em}
\sommaireligne{1}{Valeurs propres de matrices $2\times2$}{rdn-1}
\sommaireligne{2}{Retrouver une matrice à partir de ses éléments propres}{rdn-2}
\sommaireligne{3}{Une matrice triangulaire}{rdn-3}
\sommaireligne{4}{Éléments propres de $P\mapsto(X+1)P'$}{red-elem}
\sommaireligne{5}{Diagonaliser une matrice $3\times3$}{red-diag1}
\vspace{0.5em}\noindent\begin{minipage}[t]{0.86\linewidth}\raggedright\leavevmode\sffamily\bfseries\color{Ink}II. Exercices classiques et méthodes incontournables\end{minipage}\hfill\pageref{partie:II}\par\vspace{0.15em}
\sommaireligne{6}{Matrice compagnon}{red-chi}
\sommaireligne{7}{Diagonalisabilité selon un paramètre}{red-diag2}
\sommaireligne{8}{Puissances et suites couplées}{red-puiss}
\sommaireligne{9}{La matrice anti-diagonale}{rdn-4}
\sommaireligne{10}{La transposition}{rdn-5}
\sommaireligne{11}{Un endomorphisme défini avec la trace}{rdn-18}
\sommaireligne{12}{Une matrice de transition}{rdn-6}
\vspace{0.5em}\noindent\begin{minipage}[t]{0.86\linewidth}\raggedright\leavevmode\sffamily\bfseries\color{Ink}III. Exercices de synthèse et approfondissement\end{minipage}\hfill\pageref{partie:III}\par\vspace{0.15em}
\sommaireligne{13}{La matrice dont tous les coefficients valent $1$}{red-rang1}
\sommaireligne{14}{Matrices nilpotentes}{red-nilp}
\sommaireligne{15}{Matrices circulantes}{rdn-7}
\sommaireligne{16}{Un endomorphisme de $\mathbb R_n[X]$}{rdn-8}
\sommaireligne{17}{Un opérateur différentiel en dimension finie}{rdn-9}
\sommaireligne{18}{$AB$ et $BA$}{rdn-10}
\sommaireligne{19}{Une relation entre $f$ et l'identité}{rdn-11}
\sommaireligne{20}{Racines carrées d'une matrice}{red-racine}
\vspace{0.5em}\noindent\begin{minipage}[t]{0.86\linewidth}\raggedright\leavevmode\sffamily\bfseries\color{Ink}IV. Exercices difficiles et type concours\end{minipage}\hfill\pageref{partie:IV}\par\vspace{0.15em}
\sommaireligne{21}{Trigonaliser une matrice (guidé)}{red-trig}
\sommaireligne{22}{Matrices de rang 1}{rdn-12}
\sommaireligne{23}{Une récurrence linéaire d'ordre 3}{rdn-13}
\sommaireligne{24}{Racines carrées et diagonalisabilité}{rdn-14}
\end{multicols}}
\clearpage
\partie{I}{Applications directes du cours}{Valeurs propres en taille 2, lecture d'une matrice triangulaire, vecteurs propres sans matrice, une première diagonalisation.}
\exercice{1}{1}{Valeurs propres de matrices $2\times2$}{Exercice complémentaire}{polynôme caractéristique, valeurs propres complexes}{rdn-1}{1}
Déterminer les valeurs propres réelles et les sous-espaces propres de $A=\begin{pmatrix}4&-2\\1&1\end{pmatrix}$ et de
$B=\begin{pmatrix}1&-1\\1&1\end{pmatrix}$. Ces matrices sont-elles diagonalisables dans $\mathcal M_2(\R)$ ? dans $\mathcal M_2(\C)$ ?
\place
\exercice{2}{1}{Retrouver une matrice à partir de ses éléments propres}{Exercice complémentaire}{changement de base, $A=PDP^{-1}$}{rdn-2}{1}
Soit $f\in\mathcal L(\R^2)$ tel que $E_2(f)=\vect\big((1,2)\big)$ et $E_{-3}(f)=\vect\big((1,1)\big)$. Déterminer la matrice de $f$
dans la base canonique.
\place
\exercice{3}{1}{Une matrice triangulaire}{Exercice complémentaire}{matrice triangulaire, valeurs propres distinctes}{rdn-3}{2}
Montrer sans calcul que $A=\begin{pmatrix}1&0&0\\1&2&0\\1&1&3\end{pmatrix}$ est diagonalisable, puis la diagonaliser.
\place
\exercice{4}{2}{Éléments propres de $P\mapsto(X+1)P'$}{Livret, exercice 1.1}{vecteur propre, spectre, endomorphisme de polynômes}{red-elem}{2}
Soit $n\geq1$ et $f:\R_n[X]\to\R_n[X]$, $P\mapsto(X+1)P'$.
\begin{enumerate}
\item Vérifier que $f$ est un endomorphisme de $\R_n[X]$.
\item Montrer que $(X+1)^k$ est un vecteur propre de $f$ pour $0\leq k\leq n$. En déduire $\Sp(f)$ et les sous-espaces propres.
\end{enumerate}
\place
\exercice{5}{1}{Diagonaliser une matrice $3\times3$}{Livret, exercice 1.3}{polynôme caractéristique, sous-espaces propres, diagonalisation}{red-diag1}{2}
Diagonaliser $A=\begin{pmatrix}1&-1&-1\\-1&1&-1\\-1&-1&1\end{pmatrix}$ : donner $P$ inversible et $D$ diagonale telles que $A=PDP^{-1}$.
\place
\clearpage
\partie{II}{Exercices classiques et méthodes incontournables}{Polynôme caractéristique, diagonalisation selon un paramètre, puissances, symétries, matrices de transition.}
\exercice{6}{2}{Matrice compagnon}{Livret, exercice 1.2}{polynôme caractéristique, développement}{red-chi}{3}
Soient $a,b,c\in\K$ et $C=\begin{pmatrix}0&0&-c\\1&0&-b\\0&1&-a\end{pmatrix}$. Montrer que $\chi_C(x)=x^3+ax^2+bx+c$.
En déduire une matrice dont les valeurs propres sont $1$, $-1$ et $2$.
\place
\exercice{7}{2}{Diagonalisabilité selon un paramètre}{Livret, exercice 1.4}{matrice triangulaire, rang, multiplicité}{red-diag2}{3}
Pour $a\in\R$, on pose $A_a=\begin{pmatrix}1&a&1\\0&1&1\\0&0&2\end{pmatrix}$. Pour quelles valeurs de $a$ la matrice $A_a$ est-elle
diagonalisable ? La diagonaliser dans ce cas.
\place
\exercice{8}{2}{Puissances et suites couplées}{Livret, exercice 1.6}{diagonalisation, puissances, suites récurrentes}{red-puiss}{3}
Soit $A=\begin{pmatrix}1&2\\2&1\end{pmatrix}$. Calculer $A^k$ pour $k\in\N$, puis exprimer $u_k$ et $v_k$ pour les suites définies par
$u_0=1$, $v_0=0$, $u_{k+1}=u_k+2v_k$, $v_{k+1}=2u_k+v_k$, en utilisant $A^k$. Retrouver le résultat de l'exemple du cours (§1.6).
\place
\exercice{9}{2}{La matrice anti-diagonale}{Exercice complémentaire}{symétrie, sous-espaces propres, trace et déterminant}{rdn-4}{4}
Soit $n\geq2$ et $A_n\in\Mn{\R}$ définie par $a_{i,j}=1$ si $i+j=n+1$ et $a_{i,j}=0$ sinon.
\begin{enumerate}
\item Calculer $A_n^2$. En déduire que $\Sp(A_n)\subset\{-1,1\}$.
\item Déterminer $E_1(A_n)$ et $E_{-1}(A_n)$ et leurs dimensions. $A_n$ est-elle diagonalisable ?
\item En déduire $\tr A_n$ et $\det A_n$.
\end{enumerate}
\place
\exercice{10}{2}{La transposition}{Exercice complémentaire}{symétrie, matrices symétriques et antisymétriques, déterminant d'un endomorphisme}{rdn-5}{4}
Soit $T:\Mn{\R}\to\Mn{\R}$, $M\mapsto M^T$. Montrer que $T$ est diagonalisable ; déterminer ses valeurs propres, ses
sous-espaces propres, sa trace et son déterminant.
\place
\exercice{11}{2}{Un endomorphisme défini avec la trace}{Exercice complémentaire}{hyperplan, trace, diagonalisation}{rdn-18}{4}
Soit $n\geq2$ et $f:\Mn{\R}\to\Mn{\R}$, $M\mapsto M+\tr(M)I_n$.
\begin{enumerate}
\item Montrer que $H=\{M,\ \tr M=0\}$ est un hyperplan de $\Mn{\R}$ et que $\Mn{\R}=H\oplus\vect(I_n)$.
\item Montrer que $f$ est diagonalisable et calculer $\tr f$ et $\det f$. $f$ est-il un automorphisme ?
\end{enumerate}
\place
\exercice{12}{2}{Une matrice de transition}{Exercice complémentaire}{valeur propre $1$, puissances, limite}{rdn-6}{4}
Soient $p,q\in\,]0,1[$ et $A=\begin{pmatrix}1-p&q\\p&1-q\end{pmatrix}$.
\begin{enumerate}
\item Montrer que $1$ est valeur propre de $A$, et déterminer l'autre valeur propre $r$. Vérifier que $|r|<1$.
\item Diagonaliser $A$, calculer $A^k$ et sa limite quand $k\to+\infty$.
\end{enumerate}
\place
\clearpage
\partie{III}{Exercices de synthèse et approfondissement}{Matrices de rang 1, nilpotentes, circulantes ; endomorphismes de polynômes et de fonctions ; $AB$ et $BA$ ; racines carrées.}
\exercice{13}{3}{La matrice dont tous les coefficients valent $1$}{Livret, exercice 1.5}{rang, hyperplan, diagonalisation, déterminant}{red-rang1}{5}
Soit $n\geq2$ et $J_n\in\Mn{\R}$ la matrice dont tous les coefficients valent $1$.
\begin{enumerate}
\item Déterminer le rang de $J_n$, puis ses valeurs propres et ses sous-espaces propres. Est-elle diagonalisable ?
\item Soient $a,b$ réels, $b\neq0$. En déduire les valeurs propres de $aI_n+bJ_n$, et son déterminant.
\end{enumerate}
\place
\exercice{14}{2}{Matrices nilpotentes}{Livret, exercice 1.8}{valeurs propres, trace, diagonalisation}{red-nilp}{5}
Soit $A\in\Mn{\C}$ telle que $A^p=0$ pour un entier $p\geq1$.
\begin{enumerate}
\item Montrer que $\Sp(A)=\{0\}$, puis que $\chi_A(x)=x^n$. En déduire $\tr A$ et $\det A$.
\item Montrer que $A$ est diagonalisable si et seulement si $A=0$. Donner une matrice $A\neq0$ de $\mathcal M_2(\R)$ telle que $A^2=0$.
\end{enumerate}
\place
\exercice{15}{3}{Matrices circulantes}{Exercice complémentaire}{racines de l'unité, diagonalisation sur $\mathbb C$, déterminant}{rdn-7}{5}
Soit $J=\begin{pmatrix}0&0&1\\1&0&0\\0&1&0\end{pmatrix}$ et $j=e^{2i\pi/3}$.
\begin{enumerate}
\item Calculer $J^3$ et $\chi_J$. $J$ est-elle diagonalisable dans $\mathcal M_3(\R)$ ? dans $\mathcal M_3(\C)$ ?
\item Pour $\omega\in\{1,j,j^2\}$, vérifier que $X_\omega=(1,\omega^2,\omega)^T$ est vecteur propre de $J$.
\item Soient $a,b,c\in\C$ et $M=aI_3+bJ+cJ^2=\begin{pmatrix}a&c&b\\b&a&c\\c&b&a\end{pmatrix}$. Montrer que $M$ est diagonalisable et
que $\det M=(a+b+c)(a+bj+cj^2)(a+bj^2+cj)$.
\end{enumerate}
\place
\exercice{16}{3}{Un endomorphisme de $\mathbb R_n[X]$}{Exercice complémentaire}{endomorphisme de polynômes, vecteurs propres, diagonalisation}{rdn-8}{6}
Soit $n\geq1$ et $f(P)=nXP-(X^2-1)P'$ pour $P\in\R_n[X]$.
\begin{enumerate}
\item Montrer que $f$ est un endomorphisme de $\R_n[X]$.
\item Pour $0\leq k\leq n$, on pose $P_k=(X-1)^k(X+1)^{n-k}$. Montrer que $P_k$ est vecteur propre de $f$, associé à $n-2k$.
\item En déduire que $f$ est diagonalisable, et donner $\Sp(f)$, $\tr f$, $\det f$ pour $n=2$.
\end{enumerate}
\place
\exercice{17}{3}{Un opérateur différentiel en dimension finie}{Exercice complémentaire}{famille libre, noyau et image, équation différentielle}{rdn-9}{6}
Pour $k\in\ient04$, soit $e_k:x\mapsto e^{kx}$, et $E=\vect(e_0,\dots,e_4)$.
\begin{enumerate}
\item Montrer que $(e_0,\dots,e_4)$ est libre. \emph{(Ce sont des vecteurs propres de la dérivation.)}
\item Soit $\varphi(y)=y''-3y'+2y$. Montrer que $\varphi$ est un endomorphisme de $E$ ; écrire sa matrice dans la base $(e_k)$.
\item Déterminer $\Ker\varphi$ et $\Im\varphi$. Résoudre dans $E$ l'équation $\varphi(y)=e^{3x}$, et comparer avec le chapitre
\emph{Équations différentielles}.
\end{enumerate}
\place
\exercice{18}{3}{$AB$ et $BA$}{Exercice complémentaire}{valeurs propres, déterminant, diagonalisabilité}{rdn-10}{6}
Soient $A,B\in\Mn{\K}$.
\begin{enumerate}
\item Soit $\lambda\neq0$ une valeur propre de $AB$ et $X$ un vecteur propre associé. Montrer que $BX\neq0$ et que $\lambda$ est valeur
propre de $BA$.
\item Montrer que $0$ est valeur propre de $AB$ si et seulement si elle l'est de $BA$. Conclure : $\Sp(AB)=\Sp(BA)$.
\item Avec $A=\begin{pmatrix}0&1\\0&0\end{pmatrix}$ et $B=\begin{pmatrix}1&0\\0&0\end{pmatrix}$, comparer la diagonalisabilité de $AB$ et de $BA$.
\end{enumerate}
\place
\exercice{19}{3}{Une relation entre $f$ et l'identité}{Exercice complémentaire}{somme directe, noyaux, diagonalisation}{rdn-11}{7}
Soit $E$ de dimension finie et $f\in\LE$ tel que $(f-3\Id)\circ(f+2\Id)=0$.
\begin{enumerate}
\item Montrer que $\Im(f+2\Id)\subset\Ker(f-3\Id)$ et $\Im(f-3\Id)\subset\Ker(f+2\Id)$.
\item En écrivant $x=\frac15(f+2\Id)(x)-\frac15(f-3\Id)(x)$, montrer que $E=\Ker(f-3\Id)\oplus\Ker(f+2\Id)$.
\item En déduire que $f$ est diagonalisable et que $\Sp(f)\subset\{3,-2\}$.
\end{enumerate}
\place
\exercice{20}{3}{Racines carrées d'une matrice}{Livret, exercice 1.7}{commutant, matrice diagonale, changement de base}{red-racine}{7}
\begin{enumerate}
\item Déterminer les matrices $M\in\mathcal M_3(\R)$ telles que $M^2=D$, où $D=\mathrm{diag}(1,4,9)$. \emph{(Montrer d'abord que $M$ commute avec $D$.)}
\item Soit $A\in\mathcal M_3(\R)$ diagonalisable, de valeurs propres $1$, $4$, $9$. Combien existe-t-il de matrices $M$ telles que $M^2=A$ ?
\end{enumerate}
\place
\clearpage
\partie{IV}{Exercices difficiles et type concours}{Trigonalisation guidée, récurrences, matrices qui commutent, dans l'esprit des épreuves écrites.}
\exercice{21}{3}{Trigonaliser une matrice (guidé)}{Livret, exercice 1.9}{trigonalisation, multiplicité, base}{red-trig}{7}
Soit $T=\begin{pmatrix}5&-17&25\\2&-9&16\\1&-5&9\end{pmatrix}$.
\begin{enumerate}
\item Calculer $\chi_T$ et montrer que $T$ n'est pas diagonalisable, même dans $\mathcal M_3(\C)$.
\item Soient $v_1\in E_1(T)$ et $v_2\in E_2(T)$ non nuls. Trouver $v_3$ tel que $Tv_3=v_2+2v_3$ ; montrer que $(v_1,v_2,v_3)$ est une
base de $\R^3$ et donner une matrice triangulaire semblable à $T$.
\end{enumerate}
\place
\exercice{22}{4}{Matrices de rang 1}{Exercice complémentaire}{rang, trace, diagonalisation}{rdn-12}{8}
Soit $A\in\Mn{\R}$ de rang $1$ ($n\geq2$).
\begin{enumerate}
\item Montrer qu'il existe des colonnes $U,V$ non nulles telles que $A=UV^T$. Montrer que $A^2=\tr(A)\,A$.
\item Montrer que $0$ est valeur propre, avec $\dim E_0=n-1$. Montrer que $A$ est diagonalisable si et seulement si $\tr A\neq0$.
\item Exemple : $\begin{pmatrix}1&-1\\1&-1\end{pmatrix}$ est-elle diagonalisable ? et $\begin{pmatrix}1&2\\2&4\end{pmatrix}$ ?
\end{enumerate}
\place
\exercice{23}{4}{Une récurrence linéaire d'ordre 3}{Exercice complémentaire}{matrice compagnon, diagonalisation, suites}{rdn-13}{8}
Soit $(u_k)$ définie par $u_0=3$, $u_1=2$, $u_2=6$ et $u_{k+3}=2u_{k+2}+u_{k+1}-2u_k$. On pose $X_k=(u_k,u_{k+1},u_{k+2})^T$.
\begin{enumerate}
\item Écrire $X_{k+1}=CX_k$, calculer $\chi_C$ et montrer que $C$ est diagonalisable.
\item En déduire qu'il existe $\alpha,\beta,\gamma$ tels que $u_k=\alpha+\beta(-1)^k+\gamma2^k$, puis les déterminer.
\end{enumerate}
\place
\exercice{24}{4}{Racines carrées et diagonalisabilité}{Exercice complémentaire}{matrices qui commutent, sous-espaces propres stables}{rdn-14}{8}
Soient $A,B\in\mathcal M_3(\R)$ telles que $B^2=A$.
\begin{enumerate}
\item Montrer que si $B$ est diagonalisable, $A$ l'est aussi. La réciproque est-elle vraie ?
\item On suppose que $A$ a trois valeurs propres distinctes. Montrer que $AB=BA$, puis que tout vecteur propre de $A$ est vecteur
propre de $B$. En déduire que $B$ est diagonalisable.
\end{enumerate}
\place
\end{document}
